\chapter{はじめに}
\label{ch:introduction}

\section{背景}
\label{sec:intro-background}

グラフのなかで互いに全て隣接している頂点の集合を\textbf{クリーク}という。
どの頂点を足してもクリークでなくなるものを\textbf{極大クリーク}と呼び、グラフに含まれる極大クリークをすべて列挙する問題を\textbf{極大クリーク列挙}という。
極大クリークは、要素どうしが密につながった集団をそのまま取り出す構造であり、ネットワークとして表せるデータから結びつきの強いまとまりを見つける道具として使われる。

こうしたまとまりの抽出は、データ分析の一部として現れる。
扱うデータは年々大きくなり、その分析を速く回せるかどうかが、データを持つ側がそこから何を引き出せるかを左右する。
分析の一段として極大クリーク列挙を走らせるなら、その列挙自体が遅ければ後続の作業も待たされる。
本研究は、この列挙を高速化する手立てとして並列化を扱う。
たとえば、極大クリークを種として複雑ネットワークから重複しあうコミュニティを抽出する手法があり、共著ネットワークやタンパク質相互作用ネットワークに適用されている \cite{palla2005}。

\section{問題}
\label{sec:intro-problem}

極大クリーク列挙は、出力の量に敏感な問題である。
$n$ 頂点のグラフが持ちうる極大クリークの数は最悪で $3^{n/3}$ に達し、この上限は Moon と Moser の極値グラフが実際に達成する \cite{moon1965}。
出力そのものが指数的に膨らみうるため、入力サイズが同じでも、列挙にかかる時間はグラフの形によって大きく変わる。

大きなグラフでは、逐次実行の時間が実用上の壁になる。
本研究で主対象とした soc-sign-epinions（頂点数 131{,}828、極大クリーク数 約2{,}223万）では、逐次実行に、縮退順序を用いる Eppstein のアルゴリズム \cite{eppstein2010,eppstein2013} で約95秒、ピボット選択を用いる CLIQUES \cite{tomita2006} で約167秒を要した。
一度の分析で何度も列挙を回す場面を考えると、この所要時間は無視できない。
手元の計算機が複数のコアを備える以上、それらを使って列挙を速くできないかを問うのは自然である。

\section{本研究のアプローチ}
\label{sec:intro-approach}

本研究は、Eppstein のアルゴリズムの最外層を並列化する。
Eppstein のアルゴリズムは、グラフの疎さを表す\textbf{縮退数} $d$ を用いて最悪計算量を $O(d\,n\,3^{d/3})$ に抑える変種であり \cite{eppstein2010}、最外層を縮退順序に沿って回る。
この最外層の各頂点を根とする部分問題は互いに独立しているため、頂点ごとの部分問題を複数の worker へ配れば、そのまま並列化できる。

手法の選び方ではなく、進め方に本研究の重点を置く。
最外層を素朴に並列化した実装から出発し、性能を測り、伸びを止めている要因を一つ特定しては潰す、という実測駆動の反復で進めた。
律速要因を机上で見積もって一度に設計するのではなく、測定が示した順に、プロセス起動の固定費、worker 間の負荷の偏り、計算機の非対称なコア構成、並列化に伴うメモリのコピーを、順に切り分けていった。
測定は単発の値の振れが大きかったため、同一条件を3回計測して最小値を採る方針で揃えている。

\section{本研究の貢献}
\label{sec:intro-contributions}

本研究の貢献は次の四点である。

\begin{itemize}
  \item \textbf{逐次実装での二手法の実測比較}：次数分布と疎密の異なる10種類のグラフで、CLIQUES と Eppstein のアルゴリズムを同一条件で比較した。両者の優劣は縮退数のグラフサイズに対する比 $d/n$ でほぼ説明でき、疎で大きいグラフほど Eppstein が有利になる（最大で木の16.6倍）一方、密なグラフでは互角に収束し、明確に負けるのは完全グラフだけだった。この結果が、疎で大きいグラフを並列化の主対象に選ぶ根拠になっている。
  \item \textbf{律速要因の測定による切り分け}：素朴な並列実装が soc-sign-epinions で頭打ちになる原因を、フェーズ別の時間分解と頂点ごとの負荷計測で特定した。プロセスプールの起動固定費が worker 数にほぼ比例して伸びること（spawn 方式の8 worker で約1.27秒）、仕事が縮退順序の末尾の少数の重い頂点に偏ること（上位1\%の頂点が全所要時間の94.7\%を占める）、そして計算機の非対称なコア構成が効いていることを、いずれも測定値で示した。
  \item \textbf{静的な負荷分散の効果と限界}：重い頂点と軽い頂点を各バッチに混ぜる配分（interleave）で、worker の遊休を32.6\%から0.1\%までほぼ消せることを確かめた。それにもかかわらず実行時間の短縮は4.3\%にとどまり、その原因が、性能コアと効率コアからなる非対称なコア構成にあることを、コアを固定した計測で裏づけた（効率コア固定時の速度比は性能コア比0.26。これは計測方法上の下限値で、全力時は約0.43と推定される）。遊休を測る指標がコアの均質性を暗黙に仮定していたことを、この計測が明らかにした。
  \item \textbf{起動固定費の削減}：プロセス起動方式を spawn から fork に変えるだけで、worker 数に比例していた起動固定費をほぼ消せることを示した（soc-sign-epinions の8 worker で起動固定費が約0.03秒）。これにより、素朴な実装では並列化が損だった小さいグラフのうち4種が、4 worker で約2倍の高速化に転じた。共有メモリを併用する変種と合わせ、soc-sign-epinions での8 worker の実行時間を約26.5秒（逐次の94.9秒に対し約3.6倍）まで短縮した。
\end{itemize}

\section{本論文の構成}
\label{sec:intro-outline}

第2章では、クリーク、縮退数、Bron--Kerbosch 系のアルゴリズムなど、本論文で用いる用語と対象アルゴリズムを整理する。
第3章では、極大クリーク列挙の並列化に関する先行研究を、動的スケジューリングと共有メモリ化の観点から概観する \cite{schmidt2009,das2018parmce}。
第4章では、最外層並列化の素朴な実装と、測定を経て加えた改良を述べる。
第5章では、逐次比較と並列化の各段階の測定結果を示し、律速要因の切り分けを報告する。
第6章で、得られた知見と今後の課題をまとめる。
